\begin{tikzpicture}[font=\small,>=Stealth,
  box/.style={draw=Primary,fill=PrimaryLight,rounded corners=2pt,minimum width=3cm,minimum height=1.05cm,align=center,font=\scriptsize,text=PrimaryDark}]
  \node[box] (plan) at (-4.4,1.2) {泄洪方案\\孔口开度组合};
  \node[box] (gate) at (-4.4,-1.2) {闸门状态模拟\\按方案更新};
  \node[draw=Primary!40,rounded corners=3pt,minimum width=6.3cm,minimum height=3.6cm] (v) at (2.6,0) {};
  \node[anchor=north,font=\scriptsize,text=PrimaryDark] at ($(v.north)+(0,-.12)$) {闸门与泄流示意};
  \fill[Primary!20] (.05,.55)--(1.05,.55)--(1.05,.22)--(1.3,.14)--(4.8,-.83)--(4.8,-1.05)--(1.25,0)--(.05,0)--cycle;
  \draw[Primary,thick] (.05,0)--(1.25,0)--(4.8,-1.05);
  \draw[fill=TextGray!15,draw=TextGray,thick] (1.05,.3) rectangle (1.3,1.03);
  \draw[->,Primary,thick] (.2,.35)--(.85,.35);
  \draw[->,Primary,thick] (1.7,-.05)--(3.25,-.5);
  \draw[->,Primary,thick] (3.5,-.57)--(4.45,-.85);
  \node[font=\scriptsize,text=TextGray] at (3.55,.8) {开度与泄量关联同一方案};
  \node[font=\scriptsize,text=TextGray] at (2.6,-1.4) {闸下出流经溢洪道下泄};
  \draw[->,thick,Primary] (plan)--(gate);
  \draw[->,thick,Primary] (gate.east)--(v.west|-gate.east);
  \draw[->,thick,Primary] (plan.east)--(v.west|-plan.east);
\end{tikzpicture}
